% !TeX root = ../main.tex
\section{Operational semantics}
\label{sec:operational}

\subsection{The behavior of programs}

The lecture starts with a line of C++:
\begin{lstlisting}[language=C++]
foo(i++, --i);
\end{lstlisting}
\texttt{i++} returns the value of \texttt{i} and then increments it, and
\texttt{--i} decrements \texttt{i} and returns the result. With \texttt{i = 0},
what are the arguments of \texttt{foo}? If the arguments are evaluated from
left to right, \texttt{i++} gives 0 and leaves \texttt{i = 1}, then
\texttt{--i} gives 0: the call is \texttt{foo(0, 0)}. From right to left,
\texttt{--i} gives $-1$, then \texttt{i++} gives $-1$: \texttt{foo(-1, -1)}. C++
does not fix the order, and before C++17 the line was undefined. The syntax
of the line is clear; its meaning is not. Semantics is the part of the
definition of a language that gives meaning to programs.

Semantics come in three flavors: axiomatic semantics, operational semantics
(small step and big step, also called natural semantics) and denotational
semantics. Besides defining the language, they are used to prove properties
of programs, to simulate and explore their behavior and to generate
interpreters. This section follows Fern\'andez~\cite[Ch.~2, 4, 6, 8]{fernandez2014operational}.

\subsection{The language SIMP}

SIMP has commands (C), integer expressions (E) and boolean expressions (B):
\begin{lstlisting}
P   ::= C | E | B
C   ::= l := E | C ; C | if B then C else C | while B do C | skip
E   ::= !l | n | E Op E
Op  ::= + | - | * | /
B   ::= true | false | E Bop E | ¬B | B ∧ B
Bop ::= > | < | =
\end{lstlisting}
$l$ is a location of the memory (a variable), and $!l$ reads its value.
Figure~\ref{fig:simp-asts} shows the ASTs of two programs of the lecture:
the swap \texttt{z := !x ; (x := !y ; y := !z)} and the factorial
\texttt{while !l > 0 do (fac := !fac * !l ; l := !l - 1)}.

\begin{figure}[h]
\centering
\begin{forest}
  for tree={font=\small\ttfamily, s sep=3mm, l sep=4mm},
  [;
    [{:=} [z] [!x]]
    [;
      [{:=} [x] [!y]]
      [{:=} [y] [!z]]
    ]
  ]
\end{forest}
\hspace{1.2cm}
\begin{forest}
  for tree={font=\small\ttfamily, s sep=3mm, l sep=4mm},
  [while
    [> [!l] [0]]
    [;
      [{:=} [fac] [* [!fac] [!l]]]
      [{:=} [l] [- [!l] [1]]]
    ]
  ]
\end{forest}
\caption{ASTs of the swap (left) and of the factorial (right).}
\label{fig:simp-asts}
\end{figure}

In Rascal the grammar has a priority for the sequence, so that
\texttt{;} binds less than the other commands:
\coderange{Rascal}{rascal/src/Simp.rsc}{syntax}

\subsection{An abstract machine}

To give meaning to programs, the lecture first describes their behavior with
an abstract machine (Figure~\ref{fig:machine}). It has a control stack $c$
with what remains to be done, a results stack $r$ with intermediate values,
a processor for arithmetic, and a memory $m : L \to \mathbb{Z}$ from locations to
integers.

\begin{figure}[h]
\centering
\begin{tikzpicture}[box/.style={draw, minimum width=2.2cm, minimum height=8mm, font=\small, align=center}]
  \node[box] (c) at (0,0) {control\\stack $c$};
  \node[box] (r) at (3.5,0) {results\\stack $r$};
  \node[box] (alu) at (1.75,-2) {processor};
  \node[box] (m) at (6.5,-1) {memory $m$};
  \node[font=\small] (in) at (-2.6,0) {program};
  \node[font=\small] (out) at (9,-1) {result};
  \draw[flecha] (in) -- (c);
  \draw[<->, thick] (c) -- (r);
  \draw[<->, thick] (c) -- (alu);
  \draw[<->, thick] (r) -- (alu);
  \draw[<->, thick] (r) -- (m);
  \draw[flecha] (m) -- (out);
\end{tikzpicture}
\caption{The abstract machine.}
\label{fig:machine}
\end{figure}

\begin{definition}
A \textbf{configuration} of the machine is a triple $\langle c, r, m \rangle$, with
\[
c ::= \mathit{nil} \mid i \cdot c \qquad
i ::= P \mid \mathit{Op} \mid \neg \mid \wedge \mid \mathit{Bop} \mid {:=} \mid \mathit{if} \mid \mathit{while} \qquad
r ::= \mathit{nil} \mid P \cdot r \mid l \cdot r.
\]
The initial configuration for a command $C$ is $\langle C \cdot \mathit{nil}, \mathit{nil}, m \rangle$.
The machine stops at a final configuration $\langle \mathit{nil}, \mathit{nil}, m \rangle$. A
computation step transforms a configuration into a new one,
$\langle c, r, m \rangle \to \langle c', r', m' \rangle$.
\end{definition}

\paragraph{Integer expressions.}
\[
\begin{array}{lll}
\langle n \cdot c, r, m \rangle & \to \langle c, n \cdot r, m \rangle \\
\langle !l \cdot c, r, m \rangle & \to \langle c, n \cdot r, m \rangle & \text{if } m(l) = n \\
\langle E_1\ \mathit{op}\ E_2 \cdot c, r, m \rangle & \to \langle E_1 \cdot E_2 \cdot \mathit{op} \cdot c, r, m \rangle \\
\langle \mathit{op} \cdot c, n_1 \cdot n_2 \cdot r, m \rangle & \to \langle c, n \cdot r, m \rangle & \text{if } n = n_2\ \mathit{op}\ n_1
\end{array}
\]
$E_1$ is evaluated first, so its value $n_2$ ends up below the value $n_1$ of
$E_2$. The slide writes the condition as $n = n_1\ \mathit{op}\ n_2$, which gives
the right result for $+$ and $*$ but not for $-$, $/$ or the comparisons; its
own trace of the factorial computes $4 - 1 = 3$, which is $n_2 - n_1$.

Table~\ref{tab:machine-expr} is the evaluation of $3 * (4 + 5)$.

\begin{table}[h]
\centering
\small
\begin{tabular}{l l l}
  \toprule
  Control stack & Results stack & Rule \\
  \midrule
  $3 * (4 + 5) \cdot \mathit{nil}$       & $\mathit{nil}$          & split the operation \\
  $3 \cdot (4 + 5) \cdot * \cdot \mathit{nil}$ & $\mathit{nil}$     & number \\
  $(4 + 5) \cdot * \cdot \mathit{nil}$   & $3 \cdot \mathit{nil}$  & split the operation \\
  $4 \cdot 5 \cdot + \cdot * \cdot \mathit{nil}$ & $3 \cdot \mathit{nil}$ & number \\
  $5 \cdot + \cdot * \cdot \mathit{nil}$ & $4 \cdot 3 \cdot \mathit{nil}$ & number \\
  $+ \cdot * \cdot \mathit{nil}$         & $5 \cdot 4 \cdot 3 \cdot \mathit{nil}$ & $4 + 5 = 9$ \\
  $* \cdot \mathit{nil}$                 & $9 \cdot 3 \cdot \mathit{nil}$ & $3 * 9 = 27$ \\
  $\mathit{nil}$                         & $27 \cdot \mathit{nil}$ & \\
  \bottomrule
\end{tabular}
\caption{The machine on $3 * (4 + 5)$.}
\label{tab:machine-expr}
\end{table}

\paragraph{Boolean expressions.}
\[
\begin{array}{lll}
\langle b \cdot c, r, m \rangle & \to \langle c, b \cdot r, m \rangle \\
\langle \neg B \cdot c, r, m \rangle & \to \langle B \cdot \neg \cdot c, r, m \rangle \\
\langle \neg \cdot c, b \cdot r, m \rangle & \to \langle c, b' \cdot r, m \rangle & \text{if } b' = \mathit{not}\ b \\
\langle B_1 \wedge B_2 \cdot c, r, m \rangle & \to \langle B_1 \cdot B_2 \cdot \wedge \cdot c, r, m \rangle \\
\langle \wedge \cdot c, b_1 \cdot b_2 \cdot r, m \rangle & \to \langle c, b \cdot r, m \rangle & \text{if } b = b_2 \wedge b_1 \\
\langle E_1\ \mathit{Bop}\ E_2 \cdot c, r, m \rangle & \to \langle E_1 \cdot E_2 \cdot \mathit{Bop} \cdot c, r, m \rangle \\
\langle \mathit{Bop} \cdot c, n_1 \cdot n_2 \cdot r, m \rangle & \to \langle c, b \cdot r, m \rangle & \text{if } b = n_2\ \mathit{Bop}\ n_1
\end{array}
\]
(The slide writes the rule for $\neg$ as ``$b$ if $b = \mathit{not}\ b$'', which
never holds; the result is a new value $b'$.)

\paragraph{Commands.} A sequence executes from left to right. An \texttt{if}
evaluates the condition and keeps both branches on the results stack until
it knows which one to run. A \texttt{while} evaluates the condition and, if it
is true, runs the body and then the whole loop again. An assignment evaluates
the expression and then updates the memory:
\[
\begin{array}{ll}
\langle (C_1 ; C_2) \cdot c, r, m \rangle & \to \langle C_1 \cdot C_2 \cdot c, r, m \rangle \\
\langle (\mathtt{if}\ B\ \mathtt{then}\ C_1\ \mathtt{else}\ C_2) \cdot c, r, m \rangle & \to \langle B \cdot \mathit{if} \cdot c, C_1 \cdot C_2 \cdot r, m \rangle \\
\langle \mathit{if} \cdot c, \mathit{True} \cdot C_1 \cdot C_2 \cdot r, m \rangle & \to \langle C_1 \cdot c, r, m \rangle \\
\langle \mathit{if} \cdot c, \mathit{False} \cdot C_1 \cdot C_2 \cdot r, m \rangle & \to \langle C_2 \cdot c, r, m \rangle \\
\langle (\mathtt{while}\ B\ \mathtt{do}\ C) \cdot c, r, m \rangle & \to \langle B \cdot \mathit{while} \cdot c, B \cdot C \cdot r, m \rangle \\
\langle \mathit{while} \cdot c, \mathit{True} \cdot B \cdot C \cdot r, m \rangle & \to \langle C \cdot (\mathtt{while}\ B\ \mathtt{do}\ C) \cdot c, r, m \rangle \\
\langle \mathit{while} \cdot c, \mathit{False} \cdot B \cdot C \cdot r, m \rangle & \to \langle c, r, m \rangle \\
\langle (l := E) \cdot c, r, m \rangle & \to \langle E \cdot {:=} \cdot c, l \cdot r, m \rangle \\
\langle {:=} \cdot c, v \cdot l \cdot r, m \rangle & \to \langle c, r, m[l \mapsto v] \rangle
\end{array}
\]
When the assignment runs, the value of $E$ is on top of $l$; the slide writes
the stack as $l \cdot v \cdot r$, in the opposite order to its own trace.

If $\langle C \cdot \mathit{nil}, \mathit{nil}, m \rangle \to^* \langle \mathit{nil}, \mathit{nil}, m' \rangle$, the
program $C$ executed in the memory $m$ terminates successfully and produces
$m'$.

\paragraph{The factorial.} With $B = !l > 0$, $C' = \mathit{fac} := !\mathit{fac} * !l ; l := !l - 1$
and $m = \{l \mapsto 4, \mathit{fac} \mapsto 1\}$, the first iteration of
$C = \mathtt{while}\ B\ \mathtt{do}\ C'$ is in Table~\ref{tab:machine-fac}. The other three
iterations repeat the same pattern, and the machine ends in
$\langle \mathit{nil}, \mathit{nil}, m[\mathit{fac} \mapsto 24, l \mapsto 0] \rangle$.

\begin{table}[h]
\centering
\small
\begin{tabular}{l l l}
  \toprule
  Control stack & Results stack & Memory \\
  \midrule
  $C \cdot \mathit{nil}$                       & $\mathit{nil}$                     & $m$ \\
  $!l > 0 \cdot \mathit{while} \cdot \mathit{nil}$ & $B \cdot C' \cdot \mathit{nil}$ & $m$ \\
  $!l \cdot 0 \cdot {>} \cdot \mathit{while} \cdot \mathit{nil}$ & $B \cdot C' \cdot \mathit{nil}$ & $m$ \\
  $0 \cdot {>} \cdot \mathit{while} \cdot \mathit{nil}$ & $4 \cdot B \cdot C' \cdot \mathit{nil}$ & $m$ \\
  ${>} \cdot \mathit{while} \cdot \mathit{nil}$ & $0 \cdot 4 \cdot B \cdot C' \cdot \mathit{nil}$ & $m$ \\
  $\mathit{while} \cdot \mathit{nil}$           & $\mathit{True} \cdot B \cdot C' \cdot \mathit{nil}$ & $m$ \\
  $C' \cdot C \cdot \mathit{nil}$               & $\mathit{nil}$                     & $m$ \\
  $\mathit{fac} := !\mathit{fac} * !l \cdot l := !l - 1 \cdot C \cdot \mathit{nil}$ & $\mathit{nil}$ & $m$ \\
  $!\mathit{fac} * !l \cdot {:=} \cdot l := !l - 1 \cdot C \cdot \mathit{nil}$ & $\mathit{fac} \cdot \mathit{nil}$ & $m$ \\
  $!\mathit{fac} \cdot !l \cdot * \cdot {:=} \cdots$ & $\mathit{fac} \cdot \mathit{nil}$ & $m$ \\
  $!l \cdot * \cdot {:=} \cdots$                & $1 \cdot \mathit{fac} \cdot \mathit{nil}$ & $m$ \\
  $* \cdot {:=} \cdots$                         & $4 \cdot 1 \cdot \mathit{fac} \cdot \mathit{nil}$ & $m$ \\
  ${:=} \cdot l := !l - 1 \cdot C \cdot \mathit{nil}$ & $4 \cdot \mathit{fac} \cdot \mathit{nil}$ & $m$ \\
  $l := !l - 1 \cdot C \cdot \mathit{nil}$      & $\mathit{nil}$ & $m[\mathit{fac} \mapsto 4]$ \\
  $!l - 1 \cdot {:=} \cdot C \cdot \mathit{nil}$ & $l \cdot \mathit{nil}$ & $\dots$ \\
  $!l \cdot 1 \cdot - \cdot {:=} \cdot C \cdot \mathit{nil}$ & $l \cdot \mathit{nil}$ & $\dots$ \\
  $- \cdot {:=} \cdot C \cdot \mathit{nil}$     & $1 \cdot 4 \cdot l \cdot \mathit{nil}$ & $\dots$ \\
  ${:=} \cdot C \cdot \mathit{nil}$             & $3 \cdot l \cdot \mathit{nil}$ & $\dots$ \\
  $C \cdot \mathit{nil}$                        & $\mathit{nil}$ & $m[\mathit{fac} \mapsto 4, l \mapsto 3]$ \\
  \bottomrule
\end{tabular}
\caption{First iteration of the factorial on the abstract machine.}
\label{tab:machine-fac}
\end{table}

In Rascal each transition is one equation of \texttt{machineStep}, with the
head of each list as the top of the stack:
\coderange{Rascal}{rascal/src/Simp.rsc}{machine}
The tests run $3 * (4 + 5)$, the factorial and $10 - 3$ (which checks the order
of the operands).

\subsection{Small-step semantics}

The machine describes execution with stacks. Small-step semantics describes
it directly on the program, as a reduction relation that rewrites a program
into a simpler one.

\begin{definition}
A \textbf{configuration} of SIMP is a pair $\langle P, s \rangle$ of a program and a
store $s$, a map from locations to values. The program progresses by
reductions $\langle P, s \rangle \to \langle P', s' \rangle$ between configurations.
\end{definition}

The rules for expressions reduce one operand at a time, from left to right:
\begin{mathpar}
\inferrule*[right=(var)]{s(l) = n}{\langle !l, s \rangle \to \langle n, s \rangle}
\and
\inferrule*[right=(op)]{n = n_1\ \mathit{op}\ n_2}{\langle n_1\ \mathit{op}\ n_2, s \rangle \to \langle n, s \rangle}
\and
\inferrule*[right=(opL)]{\langle E_1, s \rangle \to \langle E_1', s' \rangle}{\langle E_1\ \mathit{op}\ E_2, s \rangle \to \langle E_1'\ \mathit{op}\ E_2, s' \rangle}
\and
\inferrule*[right=(opR)]{\langle E_2, s \rangle \to \langle E_2', s' \rangle}{\langle n_1\ \mathit{op}\ E_2, s \rangle \to \langle n_1\ \mathit{op}\ E_2', s' \rangle}
\end{mathpar}
The rules for booleans are the same with $\mathit{bop}$ (\textsc{bopL},
\textsc{bopR}), $\wedge$ (\textsc{andL}, \textsc{andR}) and $\neg$
(\textsc{notArg}). Writing $n_1$ in \textsc{opR} makes the order fixed: the
right operand is reduced only when the left one is already a number.

Commands:
\begin{mathpar}
\inferrule*[right={(:=R)}]{\langle E, s \rangle \to \langle E', s' \rangle}{\langle l := E, s \rangle \to \langle l := E', s' \rangle}
\and
\inferrule*[right={(:=)}]{ }{\langle l := n, s \rangle \to \langle \mathtt{skip}, s[l \mapsto n] \rangle}
\and
\inferrule*[right=(seq)]{\langle C_1, s \rangle \to \langle C_1', s' \rangle}{\langle C_1 ; C_2, s \rangle \to \langle C_1' ; C_2, s' \rangle}
\and
\inferrule*[right=(skip)]{ }{\langle \mathtt{skip} ; C, s \rangle \to \langle C, s \rangle}
\and
\inferrule*[right=(if)]{\langle B, s \rangle \to \langle B', s' \rangle}{\langle \mathtt{if}\ B\ \mathtt{then}\ C_1\ \mathtt{else}\ C_2, s \rangle \to \langle \mathtt{if}\ B'\ \mathtt{then}\ C_1\ \mathtt{else}\ C_2, s' \rangle}
\and
\inferrule*[right=(ifT)]{ }{\langle \mathtt{if}\ \mathit{True}\ \mathtt{then}\ C_1\ \mathtt{else}\ C_2, s \rangle \to \langle C_1, s \rangle}
\and
\inferrule*[right=(ifF)]{ }{\langle \mathtt{if}\ \mathit{False}\ \mathtt{then}\ C_1\ \mathtt{else}\ C_2, s \rangle \to \langle C_2, s \rangle}
\and
\inferrule*[right=(while)]{ }{\langle \mathtt{while}\ B\ \mathtt{do}\ C, s \rangle \to \langle \mathtt{if}\ B\ \mathtt{then}\ (C ; \mathtt{while}\ B\ \mathtt{do}\ C)\ \mathtt{else}\ \mathtt{skip}, s \rangle}
\end{mathpar}
\textsc{while} has no premise: it unfolds the loop once into an \texttt{if}.

\paragraph{Example.} The swap with $s(z) = 0$, $s(x) = 1$, $s(y) = 2$:
\[
\begin{array}{ll}
& \langle z := !x ; (x := !y ; y := !z), s \rangle \\
\to & \langle z := 1 ; (x := !y ; y := !z), s \rangle \\
\to & \langle \mathtt{skip} ; (x := !y ; y := !z), s[z \mapsto 1] \rangle \\
\to & \langle x := !y ; y := !z, s[z \mapsto 1] \rangle \\
\to & \langle x := 2 ; y := !z, s[z \mapsto 1] \rangle \\
\to & \langle \mathtt{skip} ; y := !z, s[z \mapsto 1, x \mapsto 2] \rangle \\
\to & \langle y := !z, s[z \mapsto 1, x \mapsto 2] \rangle \\
\to & \langle y := 1, s[z \mapsto 1, x \mapsto 2] \rangle \\
\to & \langle \mathtt{skip}, s[z \mapsto 1, x \mapsto 2, y \mapsto 1] \rangle
\end{array}
\]
At the end $s(z) = 1$, $s(x) = 2$, $s(y) = 1$. (The slides write the store of
the middle steps as $s$; it already contains the updates.)

Every step must be justified by an axiom or a rule. The first step, for
example, is the conclusion of the tree in Figure~\ref{fig:small-proof}: \textsc{seq}
reduces the first command, \textsc{:=R} reduces its expression, and \textsc{var}
reads $x$.

\begin{figure}[h]
\centering
\begin{mathpar}
\inferrule*[right=(seq)]
  {\inferrule*[right={(:=R)}]
    {\inferrule*[right=(var)]{s(x) = 1}{\langle !x, s \rangle \to \langle 1, s \rangle}}
    {\langle z := !x, s \rangle \to \langle z := 1, s \rangle}}
  {\langle z := !x ; (x := !y ; y := !z), s \rangle \to \langle z := 1 ; (x := !y ; y := !z), s \rangle}
\end{mathpar}
\caption{The derivation that justifies the first step of the swap.}
\label{fig:small-proof}
\end{figure}

\paragraph{Final configurations.} A configuration with no possible reduction is
final: $\langle n, s \rangle$ with $n$ an integer, $\langle b, s \rangle$ with $b$ a boolean,
$\langle \mathtt{skip}, s \rangle$, and $\langle !l, s \rangle$ with $l \notin \mathrm{dom}(s)$. The
last one is a \textbf{blocked} configuration. From an initial configuration
$\langle P, s \rangle$, the program executes successfully if it reaches
$\langle n, s' \rangle$ or $\langle b, s' \rangle$ (for expressions) or $\langle \mathtt{skip}, s' \rangle$
(for commands), and it fails if it reaches a blocked configuration.

In Rascal, each \texttt{step} function returns the next configuration, or an
empty list when the configuration is final:
\coderange{Rascal}{rascal/src/Simp.rsc}{small}
The tests check that the swap takes the eight steps above and ends with
$z = 1$, $x = 2$, $y = 1$, and that \texttt{x := !w} stops at a blocked
configuration.

\subsection{Big-step semantics}

Big-step semantics evaluates a program recursively, without showing each
step: a transition represents a full evaluation. $\langle P, s \rangle \Downarrow \langle v, s' \rangle$
means that $P$ in the store $s$ evaluates to $v$ and leaves $s'$.
\begin{mathpar}
\inferrule*[right=(num)]{ }{\langle n, s \rangle \Downarrow \langle n, s \rangle}
\and
\inferrule*[right=(var)]{s(l) = v}{\langle !l, s \rangle \Downarrow \langle v, s \rangle}
\and
\inferrule*[right=(op)]{\langle E_1, s \rangle \Downarrow \langle n_1, s' \rangle \\ \langle E_2, s' \rangle \Downarrow \langle n_2, s'' \rangle \\ n = n_1\ \mathit{op}\ n_2}{\langle E_1\ \mathit{op}\ E_2, s \rangle \Downarrow \langle n, s'' \rangle}
\and
\inferrule*[right={(:=)}]{\langle E, s \rangle \Downarrow \langle n, s' \rangle}{\langle l := E, s \rangle \Downarrow \langle \mathtt{skip}, s'[l \mapsto n] \rangle}
\and
\inferrule*[right=(seq)]{\langle C_1, s \rangle \Downarrow \langle \mathtt{skip}, s' \rangle \\ \langle C_2, s' \rangle \Downarrow \langle \mathtt{skip}, s'' \rangle}{\langle C_1 ; C_2, s \rangle \Downarrow \langle \mathtt{skip}, s'' \rangle}
\and
\inferrule*[right=(ifT)]{\langle B, s \rangle \Downarrow \langle \mathit{True}, s' \rangle \\ \langle C_1, s' \rangle \Downarrow \langle \mathtt{skip}, s'' \rangle}{\langle \mathtt{if}\ B\ \mathtt{then}\ C_1\ \mathtt{else}\ C_2, s \rangle \Downarrow \langle \mathtt{skip}, s'' \rangle}
\and
\inferrule*[right=(whileT)]{\langle B, s \rangle \Downarrow \langle \mathit{True}, s_1 \rangle \\ \langle C, s_1 \rangle \Downarrow \langle \mathtt{skip}, s_2 \rangle \\ \langle \mathtt{while}\ B\ \mathtt{do}\ C, s_2 \rangle \Downarrow \langle \mathtt{skip}, s_3 \rangle}{\langle \mathtt{while}\ B\ \mathtt{do}\ C, s \rangle \Downarrow \langle \mathtt{skip}, s_3 \rangle}
\and
\inferrule*[right=(whileF)]{\langle B, s \rangle \Downarrow \langle \mathit{False}, s' \rangle}{\langle \mathtt{while}\ B\ \mathtt{do}\ C, s \rangle \Downarrow \langle \mathtt{skip}, s' \rangle}
\end{mathpar}
\textsc{ifF} is like \textsc{ifT} with $\mathit{False}$ and $C_2$, and the boolean
rules are like \textsc{op}.

\paragraph{Example.} The derivation of the swap is built from the leaves up.
With $s_1 = s[z \mapsto 1]$ and $s_2 = s[z \mapsto 1, x \mapsto 2]$:
\[
\begin{array}{ll}
\langle !x, s \rangle \Downarrow \langle 1, s \rangle & \text{(var)} \\
\langle z := !x, s \rangle \Downarrow \langle \mathtt{skip}, s_1 \rangle & \text{(:=)} \\
\langle !y, s_1 \rangle \Downarrow \langle 2, s_1 \rangle & \text{(var)} \\
\langle x := !y, s_1 \rangle \Downarrow \langle \mathtt{skip}, s_2 \rangle & \text{(:=)} \\
\langle !z, s_2 \rangle \Downarrow \langle 1, s_2 \rangle & \text{(var)} \\
\langle y := !z, s_2 \rangle \Downarrow \langle \mathtt{skip}, s_2[y \mapsto 1] \rangle & \text{(:=)} \\
\langle z := !x ; x := !y ; y := !z, s \rangle \Downarrow \langle \mathtt{skip}, s[z \mapsto 1, x \mapsto 2, y \mapsto 1] \rangle & \text{(seq), twice}
\end{array}
\]
Figure~\ref{fig:big-proof} draws the tree for the first two assignments.

\begin{figure}[h]
\centering
\begin{mathpar}
\inferrule*[right=(seq)]
  {\inferrule*[right={(:=)}]{\inferrule*[right=(var)]{ }{\langle !x, s \rangle \Downarrow \langle 1, s \rangle}}{\langle z := !x, s \rangle \Downarrow \langle \mathtt{skip}, s_1 \rangle}
   \\
   \inferrule*[right={(:=)}]{\inferrule*[right=(var)]{ }{\langle !y, s_1 \rangle \Downarrow \langle 2, s_1 \rangle}}{\langle x := !y, s_1 \rangle \Downarrow \langle \mathtt{skip}, s_2 \rangle}}
  {\langle z := !x ; x := !y, s \rangle \Downarrow \langle \mathtt{skip}, s_2 \rangle}
\end{mathpar}
\caption{Big-step derivation of the first two assignments of the swap.}
\label{fig:big-proof}
\end{figure}

In Rascal big-step semantics is a recursive evaluator with one equation per
rule:
\coderange{Rascal}{rascal/src/Simp.rsc}{big}
Small step is the natural model of a debugger: it shows each intermediate
state, and it can describe programs that do not terminate (an infinite
sequence of steps). Big step is the natural model of an interpreter, and a
program that does not terminate simply has no derivation.

\subsection{Blocks and local variables}

Blocks give locality to programs: a variable declared in a block is only
visible and valid inside it. SIMP is extended with blocks and with
allocation of new locations:
\begin{lstlisting}
C ::= ... | begin loc l := E ; C end
E ::= ... | ref E
\end{lstlisting}
\begin{mathpar}
\inferrule*[right=(ref)]{\langle E, s \rangle \Downarrow \langle v, s' \rangle}{\langle \mathtt{ref}\ E, s \rangle \Downarrow \langle l, s'[l \mapsto v] \rangle}\quad l \notin \mathrm{dom}(s')
\and
\inferrule*[right=(loc)]{\langle E, s \rangle \Downarrow \langle v, s' \rangle \\ \langle C\{z \mapsto l\}, s'[l \mapsto v] \rangle \Downarrow \langle \mathtt{skip}, s''[l \mapsto v'] \rangle}{\langle \mathtt{begin}\ \mathtt{loc}\ z := E ; C\ \mathtt{end}, s \rangle \Downarrow \langle \mathtt{skip}, s'' \rangle}
\end{mathpar}
with $l \notin \mathrm{dom}(s') \cup \mathrm{dom}(s'') \cup \mathit{locations}(C)$. $C\{z \mapsto l\}$
replaces the name $z$ by the new location $l$ in $C$, and the final store
drops $l$: the local variable does not exist after the block. For the swap
with a local variable,
\begin{lstlisting}
begin
  loc z := !x;
  x := !y;
  y := !z
end
\end{lstlisting}
the body runs in $s'[l \mapsto s(x)]$ and gives $s_1 = s[l \mapsto s(x), x \mapsto s(y), y \mapsto s(x)]$;
the whole program gives $s_2 = s[x \mapsto s(y), y \mapsto s(x)]$, with no $z$.
\coderange{Rascal}{rascal/src/Simp.rsc}{blockbig}

\subsection{More constructs}

This subsection goes beyond the slides.

\paragraph{repeat.} The command $\mathtt{repeat}\ E\ \mathtt{times}\ C$ runs $C$ as many
times as the value of $E$, evaluated once. In small step:
\begin{mathpar}
\inferrule*[right=(repeat-E)]{\langle E, s \rangle \to \langle E', s' \rangle}{\langle \mathtt{repeat}\ E\ \mathtt{times}\ C, s \rangle \to \langle \mathtt{repeat}\ E'\ \mathtt{times}\ C, s' \rangle}
\and
\inferrule*[right=(repeat-fin)]{n \le 0}{\langle \mathtt{repeat}\ n\ \mathtt{times}\ C, s \rangle \to \langle \mathtt{skip}, s \rangle}
\and
\inferrule*[right=(repeat-paso)]{n > 0 \\ m = n - 1}{\langle \mathtt{repeat}\ n\ \mathtt{times}\ C, s \rangle \to \langle C ; \mathtt{repeat}\ m\ \mathtt{times}\ C, s \rangle}
\end{mathpar}
Once $E$ is a numeral it is never evaluated again, even if $C$ changes the
variables of $E$. The base case is $n \le 0$ and not $n = 0$: with only
$n = 0$, a negative count would leave the configuration blocked.
\coderange{Rascal}{rascal/src/Simp.rsc}{repeatsmall}
\coderange{Rascal}{rascal/src/Simp.rsc}{repeatbig}
The tests run \texttt{repeat !k times (x := !x + 1; k := !k + 1)} with $k = 2$,
which runs the body twice although $k$ grows, and a negative count.

\paragraph{for.} A loop $\mathtt{for}\,(C_i; B; C_s)\,\{C_b\}$ is syntactic sugar for
$C_i ; \mathtt{while}\ B\ \mathtt{do}\ (C_b ; C_s)$. With its own rules, the
initialization is reduced first and then the loop unfolds into an
\texttt{if}:
\begin{mathpar}
\inferrule*[right=For-Init]{\langle C_i, s \rangle \to \langle C_i', s' \rangle}{\langle \mathtt{for}\,(C_i; B; C_s)\,\{C_b\}, s \rangle \to \langle \mathtt{for}\,(C_i'; B; C_s)\,\{C_b\}, s' \rangle}
\and
\inferrule*[right=For-Skip]{ }{\langle \mathtt{for}\,(\mathtt{skip}; B; C_s)\,\{C_b\}, s \rangle \to \langle \mathtt{if}\ B\ \mathtt{then}\ (C_b ; C_s ; \mathtt{for}\,(\mathtt{skip}; B; C_s)\,\{C_b\})\ \mathtt{else}\ \mathtt{skip}, s \rangle}
\end{mathpar}
\coderange{Rascal}{rascal/src/Simp.rsc}{forsmall}

\paragraph{Lists.} Big step also defines functions on data. With the
constructors $\nil$ and $\cons{h}{t}$, adding an element at the end of a list
and appending two lists are
\begin{mathpar}
\inferrule*[right=Concat-Nil]{ }{\bigstep{\concat{e}{\nil}}{\sigma}{\cons{e}{\nil}}}
\and
\inferrule*[right=Concat-Cons]{\bigstep{\concat{e}{t}}{\sigma}{t'}}{\bigstep{\concat{e}{\cons{h}{t}}}{\sigma}{\cons{h}{t'}}}
\and
\inferrule*[right=Append-Nil]{ }{\bigstep{\append{\nil}{\ell_2}}{\sigma}{\ell_2}}
\and
\inferrule*[right=Append-Cons]{\bigstep{\append{t_1}{\ell_2}}{\sigma}{\ell'}}{\bigstep{\append{\cons{h}{t_1}}{\ell_2}}{\sigma}{\cons{h}{\ell'}}}
\end{mathpar}

\subsection{Exercises}

\begin{exercise}
Draw the AST of \texttt{while !l > 0 do (factorial := !factorial * !l; l := !l - 1)}.
\end{exercise}

\paragraph{Solution.} It is the right tree of Figure~\ref{fig:simp-asts}, with
\texttt{factorial} in place of \texttt{fac}.

\begin{exercise}
Define the transitions of the abstract machine for \texttt{;}, \texttt{if},
\texttt{while} and assignment.
\end{exercise}

\paragraph{Solution.} They are the command rules of the abstract machine above.
The idea in each one: put the parts that must run first on the control stack,
and keep on the results stack what will be needed later (the branches of the
\texttt{if}, the condition and the body of the \texttt{while}, the location of
the assignment).

\begin{exercise}
Run \texttt{y := !x + 1 ; x := !y * 2} from $\sigma = \{x \mapsto 3, y \mapsto 0\}$ in small
step and in big step.
\end{exercise}

\paragraph{Solution.} In small step, with $\sigma_1 = \sigma[y \mapsto 4]$ and
$\sigma_2 = \sigma_1[x \mapsto 8]$:
\[
\begin{array}{ll}
& \langle y := !x + 1 ; x := !y * 2, \sigma \rangle \\
\to & \langle y := 3 + 1 ; x := !y * 2, \sigma \rangle \quad \text{(seq, :=R, opL, var)} \\
\to & \langle y := 4 ; x := !y * 2, \sigma \rangle \quad \text{(seq, :=R, op)} \\
\to & \langle \mathtt{skip} ; x := !y * 2, \sigma_1 \rangle \quad \text{(seq, :=)} \\
\to & \langle x := !y * 2, \sigma_1 \rangle \quad \text{(skip)} \\
\to & \langle x := 4 * 2, \sigma_1 \rangle \to \langle x := 8, \sigma_1 \rangle \to \langle \mathtt{skip}, \sigma_2 \rangle
\end{array}
\]
In big step there is one tree: \textsc{seq} with two premises, each an
assignment whose premise evaluates $!x + 1 \Downarrow 4$ in $\sigma$ and $!y * 2 \Downarrow 8$
in $\sigma_1$. Both give $\{x \mapsto 8, y \mapsto 4\}$.

\begin{exercise}
Give small-step rules for $\mathtt{for}\ E_1\ \mathtt{to}\ E_2\ \mathtt{do}\ C$, which evaluates
$E_1$ and $E_2$ once and runs $C$ for each integer from $n_1$ to $n_2$.
\end{exercise}

\paragraph{Solution.} Reduce the bounds first, then unfold:
\begin{mathpar}
\inferrule{\langle E_1, s \rangle \to \langle E_1', s' \rangle}{\langle \mathtt{for}\ E_1\ \mathtt{to}\ E_2\ \mathtt{do}\ C, s \rangle \to \langle \mathtt{for}\ E_1'\ \mathtt{to}\ E_2\ \mathtt{do}\ C, s' \rangle}
\and
\inferrule{\langle E_2, s \rangle \to \langle E_2', s' \rangle}{\langle \mathtt{for}\ n_1\ \mathtt{to}\ E_2\ \mathtt{do}\ C, s \rangle \to \langle \mathtt{for}\ n_1\ \mathtt{to}\ E_2'\ \mathtt{do}\ C, s' \rangle}
\and
\inferrule{n_1 > n_2}{\langle \mathtt{for}\ n_1\ \mathtt{to}\ n_2\ \mathtt{do}\ C, s \rangle \to \langle \mathtt{skip}, s \rangle}
\and
\inferrule{n_1 \le n_2}{\langle \mathtt{for}\ n_1\ \mathtt{to}\ n_2\ \mathtt{do}\ C, s \rangle \to \langle C ; \mathtt{for}\ (n_1 + 1)\ \mathtt{to}\ n_2\ \mathtt{do}\ C, s \rangle}
\end{mathpar}
Unfolding before the bounds are numbers would evaluate them again in each
iteration. Without the rule for $n_1 > n_2$ an empty range would be blocked.

\begin{exercise}
Write big-step rules for $\mathbf{reverse}(\ell)$, using \texttt{concat}, and prove
\[
\bigstep{\concat{3}{\cons{1}{\cons{2}{\nil}}}}{\sigma}{\cons{1}{\cons{2}{\cons{3}{\nil}}}}.
\]
\end{exercise}

\paragraph{Solution.}
\begin{mathpar}
\inferrule*[right=Rev-Nil]{ }{\mathbf{reverse}(\nil) \Downarrow \nil}
\and
\inferrule*[right=Rev-Cons]{\mathbf{reverse}(t) \Downarrow t' \\ \concat{h}{t'} \Downarrow \ell'}{\mathbf{reverse}(\cons{h}{t}) \Downarrow \ell'}
\end{mathpar}
For the
proof, \textsc{Concat-Cons} twice and \textsc{Concat-Nil} once:
$\concat{3}{\nil} \Downarrow \cons{3}{\nil}$, then
$\concat{3}{\cons{2}{\nil}} \Downarrow \cons{2}{\cons{3}{\nil}}$, then the goal.
